
\exo

\begin{enumerate}
\item Parmi les nombres suivants, lesquels sont décimaux ?
\[\frac{3}{10} \qpvq 101 \qpvq \frac{1}{6} \qpvq 0,077 \qpvq -\frac{1}{5} \qpvq \frac{\pi}{10} \qpvq -13.\]
\item Parmi les nombres suivants, lesquels sont rationnels ?
\[\frac{1}{10} \qpvq -\frac{2}{3} \qpvq \frac{1}{\sqrt{2}} \qpvq 10 \qpvq -14,58.\]
\end{enumerate}

\exo

\begin{enumerate}
\item Justifier que $\frac{2}{5}$ et $\frac{3}{2}$ sont des nombres décimaux, puis déterminer leurs inverses.
\item L'inverse d'un nombre décimal (non nul) est-il encore un nombre décimal ?
\end{enumerate}

\exo

Soient deux nombres décimaux $x$ et $y$.
\begin{enumerate}
\item Montrer que le produit $x\times y$ est un nombre décimal.
\item La somme $x+y$ est-elle un nombre décimal ?
\end{enumerate}

\exo

Soient les deux nombres $a=\frac{3}{4}$ et $b=\frac{4}{5}$.
\begin{enumerate}
\item Montrer que $a$ et $b$ sont des nombres décimaux.
\item Comparer $a$ et $b$.
\item Déterminer un nombre décimal strictement compris entre $a$ et $b$.
\end{enumerate}

\exo

\begin{enumerate}
\item Justifier que les nombres $\frac{1}{4}$, $\frac{3}{25}$, $\frac{7}{20}$, $\frac{11}{2^3\times 5^2}$ sont des décimaux.
\item Soient $a$ un entier relatif, $p$ et $q$ des entiers naturels.

Montrer que le nombre $\dfrac{a}{2^p\times 5^q}$ est un décimal.
\end{enumerate}

\exo

Écrire les nombres suivants sous forme d'une fraction irréductible :
\[
a=\frac{76}{24}
\qpvq b=\frac{225}{135}
\qpvq c=\frac{180}{108}
\qpvq d=\frac{143}{52}
\qpvq e=\frac{420}{2520}
\qpvq f=\frac{110}{225}
\qpvq g=\frac{168}{180}.
\]

\exo

Vrai ou faux ? Justifier la réponse.
\begin{enumerate}
\item Le quotient de deux nombres décimaux non nuls est un nombre décimal.
\item Le produit d'un nombre décimal par $10$ est un nombre entier.
\item Le quotient de deux nombres rationnels non nuls est un nombre rationnel.
\item L'inverse d'un nombre entier non nul est un nombre décimal.
\end{enumerate}

\exo

Pour chacun des nombres suivants, dire auxquels des ensembles $\N$, $\Z$, $\D$, $\Q$ et $\R$ il appartient.

\vspace{-6pt}
\setlength{\columnsep}{1cm}
\setlength{\columnseprule}{0.1pt}
\begin{multicols}{4}
\begin{enumerate}
\item $a=0,777\vphantom{\dfrac{\sqrt{25}}{\sqrt{9}}}$

\medskip
$b=\dfrac{7}{25}\vphantom{\dfrac{\sqrt{25}}{\sqrt{9}}}$

\medskip
$c=3-\sqrt{36}\vphantom{\dfrac{\sqrt{25}}{\sqrt{9}}}$
\item $a=\dfrac{1}{6}\vphantom{\dfrac{\sqrt{25}}{\sqrt{9}}}$

\medskip 
$b=\dfrac{\pi}{4}\vphantom{\dfrac{\sqrt{25}}{\sqrt{9}}}$

\medskip
$c=10^{-4}\vphantom{\dfrac{\sqrt{25}}{\sqrt{9}}}$
\item $a=-14,6\vphantom{\dfrac{\sqrt{25}}{\sqrt{9}}}$

\medskip 
$b=\dfrac{197}{2}\vphantom{\dfrac{\sqrt{25}}{\sqrt{9}}}$

\medskip
$c=10^{10}\vphantom{\dfrac{\sqrt{25}}{\sqrt{9}}}$
\item $a=3,14159\vphantom{\dfrac{\sqrt{25}}{\sqrt{9}}}$

\medskip
$b=\sqrt{\pi^2}\vphantom{\dfrac{\sqrt{25}}{\sqrt{9}}}$

\medskip
$c=\dfrac{\sqrt{25}}{\sqrt{9}}$
\end{enumerate}
\end{multicols}
\vspace{-12pt}

\exo

\begin{enumerate}
\item Le nombre $0$ est-il irrationnel ?
\item En raisonnant par l'absurde, démontrer que l'inverse d'un nombre irrationnel est un nombre irrationnel.
\end{enumerate}

\exo

On a vu en cours que $\sqrt{2}$ est un nombre irrationnel.

En utilisant ce résultat et en raisonnant par l'absurde, démontrer que :
\begin{enumerate}
\item le nombre $a=5\sqrt{2}$ est un nombre irrationnel ;
\item le nombre $b=3+\sqrt{2}$ est un nombre irrationnel ;
\item le nombre $c=\sqrt{\sqrt{2}}$ est un nombre irrationnel.
\end{enumerate}

\exo

Répondre aux questions suivantes en utilisant la calculatrice, et et en arrondissant le résultat avec un nombre de chiffres significatifs adapté.
\begin{enumerate}
\item Une pièce contient $n=\SI{242}{\mole}$ de dioxygène, de masse molaire $M=\SI{31,9988}{\gram\per\mole}$.

Quelle est la masse de dioxygène dans la pièce ?

\medskip
\item Aux bornes d'un dipôle de résistance $R=\SI{330}{\ohm}$, on a mesuré une tension $U=\SI{2,789}{\volt}$.

Quelle est l'intensité $I$ du courant qui parcourt le dipôle ? On rappelle la loi d'Ohm : $U=R\times I$.

\medskip
\item La masse volumique d'un corps de masse $m$ (en kg) et de volume $V$ (en m\up{3}) est donnée par la formule $\rho=\frac{m}{V}$, en \si{\kilo\gram\per\m^3}.

Calculer la masse volumique de la Terre, sachant que son rayon vaut \np{6400} kilomètres, et que sa masse vaut $6\times 10^{24}$ kilogrammes.
\end{enumerate}

\exo

Dans chaque cas, donner un encadrement d'amplitude $10^{-n}$ du nombre $x$.

\vspace{-6pt}
\setlength{\columnsep}{1cm}
\setlength{\columnseprule}{0pt}
\begin{multicols}{3}
\begin{enumerate}
\item
\begin{enumerate}
\item $x=\sqrt{11}$ et $n=5$.

\medskip
\item $x=\dfrac{113}{31}$ et $n=7$.
\end{enumerate}
\item 
\begin{enumerate}
\item $x=\dfrac{1}{17}$ et $n=1$.

\medskip
\item $x=\dfrac{\sqrt{5}}{7}$ et $n=4$.
\end{enumerate}
\item
\begin{enumerate}
\item $x=\dfrac{89}{173}$ et $n=5$.

\medskip
\item $x=\sqrt{5-\sqrt{5}}$ et $n=3$.
\end{enumerate}
\end{enumerate}
\end{multicols}

\exo

Sur la droite graduée ci-dessous, on a représenté un ensemble de réels noté $I$ :
\begin{center}
\def\xmin{-2.5}
\def\xmax{3.5}
\def\ymin{-0.4}
\def\ymax{0.4}
\def\xunite{1.5cm}
\def\yunite{1.5cm}
\psset{unit=\xunite, xunit=\xunite, yunit=\yunite, algebraic=true, dimen=middle, dotstyle=*, dotsize=3pt 0, linewidth=0.8pt, arrowsize=3pt 2, arrowinset=0.25, comma=true}
\begin{pspicture*}(\xmin,\ymin)(\xmax,\ymax)
\psgrid[xunit=\xunite, yunit=\yunite, subgriddiv=0, gridlabels=0, gridcolor=lightgray](0,0)(\xmin,\ymin)(\xmax,\ymax)
\psaxes[labelFontSize=\scriptstyle, xAxis=true, yAxis=false, Dx=1, Dy=1, ticksize=-2pt 0, subticks=1]{->}(0,0)(\xmin,\ymin)(\xmax,\ymax)
\psline[linewidth=2.0pt]{[-[}(-1,0)(2,0)
\begin{footnotesize}
\uput{3pt}[090](1.2,0){$I$}
\end{footnotesize}
\end{pspicture*}
\end{center}
\begin{enumerate}
\item Écrire cet ensemble sous la forme d'un intervalle.
\item Parmi les nombres suivants, préciser ceux qui appartiennent à $I$ :
\[-0,5
\qpvq 1,999
\qpvq
2
\qpvq
\frac{5}{3}.\]
\end{enumerate}

\exo

Représenter sur la droite numérique les intervalles :
\[
I=\intof{-\infty}{1}
\qpvq
J=\intoo{-2}{+\infty}
\qpvq
K=\intfo{3}{5}.
\]

\exo

Compléter les phrases suivantes en utilisant le symbole d'appartenance ($\in$), ou le symbole de non-appartenance ($\notin$) :

\vspace{-6pt}
\setlength{\columnsep}{1cm}
\setlength{\columnseprule}{0pt}
\begin{multicols}{3}
\begin{enumerate}
\item $3\quad\dots{}\quad\intof{-1}{8}$

\medskip
\item $10^{-3}\quad\dots{}\quad\intfo{0}{+\infty}$
\item $-2\quad\dots{}\quad\intoo{-\infty}{-2}$

\medskip
\item $-2\quad\dots{}\quad\intof{-1}{6}$
\item $\pi\quad\dots{}\quad\intoo{3,14}{3,15}$

\medskip
\item $0\quad\dots{}\quad\intfo{-\sqrt{2}}{\sqrt{2}}$
\end{enumerate}
\end{multicols}

\exo

Dans chaque cas, représenter sur un axe l'ensemble donné, puis l'écrire sous la forme d'un intervalle :

\begin{enumerate}
\item $I$ est l'ensemble des réels supérieurs ou égaux à $3$.
\item $J$ est l'ensemble des réels strictement compris entre $-2$ et $4$.
\item $K$ est l'ensemble des réels strictement inférieurs à $1$.
\item $L$ est l'ensemble des réels positifs ou nuls, et strictement inférieurs à $5$.
\end{enumerate}

\exo

Dans chaque cas, traduire la condition sur $x$ sous la forme de l'appartenance de $x$ à un intervalle :

\vspace{-6pt}
\setlength{\columnsep}{1cm}
\setlength{\columnseprule}{0pt}
\begin{multicols}{3}
\begin{enumerate}
\item $-5\leqslant x\leqslant 1$

\medskip
\item $x<4$

\medskip
\item $x>7$

\medskip
\item $1<x<2$

\medskip
\item $x\geqslant 0,5$

\medskip
\item $x\leqslant 10^{-1}$

\medskip
\item $\frac{1}{2}\leqslant x< \frac{3}{4}$

\medskip
\item $x>\frac{\pi}{2}$

\medskip
\item $x<-3$
\end{enumerate}
\end{multicols}

\exo

Dans chaque cas, traduire la condition sur $x$ par des inégalités :

\vspace{-6pt}
\setlength{\columnsep}{1cm}
\setlength{\columnseprule}{0pt}
\begin{multicols}{3}
\begin{enumerate}
\item $x\in\intof{-3}{3}$

\medskip
\item $x\in\intff{5}{9}$

\medskip
\item $x\in\intoo{-0,5}{+\infty}$

\medskip
\item $x\in\intof{-7}{0}$

\medskip
\item $x\in\intoo{-\infty}{\pi}$

\medskip
\item $x\in\intfo{-1}{10}$

\medskip
\item $x\in\intfo{-1}{+\infty}$

\medskip
\item $x\in\intoo{-3}{3}$

\medskip
\item $x\in\intoo{0}{+\infty}$
\end{enumerate}
\end{multicols}

\exo

\begin{enumerate}
\item Déterminer l'amplitude de l'intervalle $\intoo{-2,5}{3,8}$.
\item Déterminer la valeur du réel $a$, de sorte que l'amplitude de l'intervalle $\intff{\dfrac{2}{3}}{a}$ soit égale à $\dfrac{1}{10}$.
\item Déterminer la valeur du réel $b$, de sorte que l'amplitude de l'intervalle $\intff{0,8-\dfrac{1}{b}}{0,8+\dfrac{1}{b}}$ soit égale à~$0,04$.
\end{enumerate}

\exo

Dans chaque cas, déterminer les ensembles $I\cap J$ et $I\cup J$ :
\vspace{-6pt}
\setlength{\columnsep}{1cm}
\setlength{\columnseprule}{0pt}
\begin{multicols}{2}
\begin{enumerate}
\item $I=\intfo{-2,5}{4}$ et $J=\intof{1}{5}$

\medskip
\item $I=\intoo{-\infty}{4}$ et $J=\intfo{2}{+\infty}$

\medskip
\item $I=\intfo{-1}{+\infty}$ et $J=\intof{-2}{3}$

\medskip
\item $I=\intff{2}{5,5}$ et $J=\intoo{1}{2}$
\end{enumerate}
\end{multicols}

\exo

Soit l'intervalle $I=\intfo{-5}{3}$.

Déterminer les ensembles $I\cap\Z$ et $I\cap\N$.

\exo

On note $x_A$ et $x_B$ les abscisses respectives de deux points $A$ et $B$ d'une droite graduée.

Dans chaque cas, déterminer la distance $AB$.

\vspace{-6pt}
\setlength{\columnsep}{1cm}
\setlength{\columnseprule}{0pt}
\begin{multicols}{3}
\begin{enumerate}
\item $x_A=-1$ et $x_B=5$.
\item $x_A=15$ et $x_B=7$.
\item $x_A=\frac{7}{4}$ et $x_B=-\frac{3}{2}$.
\end{enumerate}
\end{multicols}

\exo

Dans chaque cas, interpréter l'égalité donnée en termes de distance, puis déterminer tous les nombres réels $x$ qui la vérifient.

\vspace{-6pt}
\setlength{\columnsep}{1cm}
\setlength{\columnseprule}{0pt}
\begin{multicols}{3}
\begin{enumerate}
\item 
\begin{enumerate}
\item $\abs{x-2}=1$.

\medskip
\item $\abs{x-5}=2$.
\end{enumerate}
\item 
\begin{enumerate}
\item $\abs{x+4}=5$.

\medskip
\item $\abs{x}=2$.
\end{enumerate}
\item 
\begin{enumerate}
\item $\abs{x+2}=0$.

\medskip
\item $\abs{x-3}=3$.
\end{enumerate}
\end{enumerate}
\end{multicols}

\exo

Dans chaque cas, représenter l'intervalle $\intff{a-r}{a+r}$ pour les valeurs données de $a$ et $r$, puis caractériser l'appartenance à cet intervalle par une inégalité faisant intervenir une valeur absolue.

\vspace{-6pt}
\setlength{\columnsep}{1cm}
\setlength{\columnseprule}{0pt}
\begin{multicols}{3}
\begin{enumerate}
\item 
\begin{enumerate}
\item $a=1$ et $r=1$.

\medskip
\item $a=-3$ et $r=2$.
\end{enumerate}
\item 
\begin{enumerate}
\item $a=10$ et $r=5$.

\medskip
\item $a=-7$ et $r=0,5$.
\end{enumerate}
\item 
\begin{enumerate}
\item $a=0$ et $r=\frac{1}{2}$.

\medskip
\item $a=\frac{1}{2}$ et $r=2$.
\end{enumerate}
\end{enumerate}
\end{multicols}

\exo

Dans chaque cas, déterminer l'ensemble des valeurs du réel $x$ vérifiant l'inégalité donnée.

\vspace{-6pt}
\setlength{\columnsep}{1cm}
\setlength{\columnseprule}{0pt}
\begin{multicols}{3}
\begin{enumerate}
\item 
\begin{enumerate}
\item $\abs{x-2}\leqslant 1$.

\medskip
\item $\abs{x-7}\leqslant 4$.
\end{enumerate}
\item 
\begin{enumerate}
\item $\abs{x+3}\leqslant 1$.

\medskip
\item $\abs{x+5}\leqslant 2$.
\end{enumerate}
\item 
\begin{enumerate}
\item $\abs{x-\frac{1}{2}}\leqslant \frac{1}{4}$.

\medskip
\item $\abs{x-1}< 10^{-2}$.
\end{enumerate}
\end{enumerate}
\end{multicols}

\exo

\begin{enumerate}
\item Déterminer et représenter l'ensemble $A$ des réels $x$ tels que $\abs{x-8}\leqslant 4$.
\item En déduire  l'ensemble $B$ des réels $x$ tels que $\abs{x-8}> 4$ (on écrira $B$ en utilisant des intervalles).
\item Que représente l'ensemble $B$ pour l'ensemble $A$ ?
\end{enumerate}

\exo

On considère les deux nombres :
\[
A=\frac{\np{777777777777775}}{\np{777777777777774}}
\qetq
B=\frac{\np{777777777777774}}{\np{777777777777775}}.
\]
\begin{enumerate}
\item Classer les trois nombres $A$, $B$ et $1$.
\item Calculer $C=A-1$ et $D=1-B$.
\item Comparer $C$ et $D$.
\item Parmi les deux nombres $A$ et $B$, quel est celui qui est le plus proche de $1$ ?
\end{enumerate}

